\documentclass[11pt]{article}
\usepackage[margin=2cm]{geometry}
\usepackage{booktabs}
\usepackage[table]{xcolor}
\usepackage{amsmath}
\usepackage{amssymb}
\usepackage{parskip}

\title{Result Tables}
\author{Roham izadidoost}
\date{21/aug}

\begin{document}
\maketitle

All numbers below are means over \textbf{10 seeds} (previous version of this
document reported 3--5 seeds). Seed count was raised because Wilcoxon
signed-rank at $n=5$ cannot reach $p<0.05$ at all: with every difference the
same sign the minimum attainable two-sided $p$ is $2/2^5 = 0.0625$. That is a
combinatorial floor, not a power problem.

\section*{Main results}
Each column is a held-out target, each row a different method
(Source-only = no adaptation at all. BN-only = no-op control, since our encoder
uses LayerNorm and has no BatchNorm statistics to recalibrate. Tent = a known
TTA baseline. Self-training only = ours without the consistency term. Ours =
full method. Oracle = supervised upper bound). Ours wins on 3 targets out of 4;
each cell is EER\%\,/\,AUC, lower/higher is better.
\begin{table}[h]
\centering
\begin{tabular}{@{}lcccc@{}}
\toprule
Method & ASVspoof2019 & LibriSpeech-TTS & In-the-Wild & Arabic \\
\midrule
Source-only & 5.03 / .989 & 33.54 / .727 & 12.78 / .945 & 22.50 / .843 \\
BN-only (no-op) & 5.03 / .989 & 33.54 / .727 & 12.78 / .945 & 22.50 / .843 \\
Tent & 23.47 / .776 & 35.44 / .674 & 27.34 / .735 & 47.35 / .527 \\
Self-training only & 4.20 / .987 & \textbf{33.39} / .709 & 11.55 / .950 & 22.19 / .846 \\
\rowcolor{green!12}\textbf{Ours (full TTA)} & \textbf{3.47} / .994 & 33.74 / .727 & \textbf{11.33} / .958 & \textbf{20.99} / .866 \\
\midrule
Oracle (supervised upper bound) & 0.89 / .999 & 13.93 / .939 & 4.42 / .992 & 9.50 / .968 \\
\bottomrule
\end{tabular}
\end{table}

BN-only reproduces Source-only to every reported digit on all four targets.
That is the expected result, not a bug: it is the architectural argument for
why adaptation here has to be gradient-based.

Tent is much worse at 10 seeds than it looked at 3 --- it now degrades every
target, catastrophically on Arabic (47.35\%, i.e.\ chance).

\section*{Seed variability (bootstrap 95\% CIs)}
Bootstrap confidence intervals over the 10 seeds. Source and Ours are
non-overlapping on 3 of 4 targets; they overlap exactly on LibriSpeech-TTS,
which is the target where we do not claim a gain.
\begin{table}[h]
\centering
\begin{tabular}{@{}lccc@{}}
\toprule
Target & Source-only EER (95\% CI) & Ours EER (95\% CI) & Disjoint? \\
\midrule
ASVspoof2019 & 5.03 [4.27, 5.82] & \textbf{3.47} [2.97, 3.99] & yes \\
In-the-Wild & 12.78 [12.06, 13.50] & \textbf{11.33} [10.53, 12.21] & yes \\
Arabic & 22.50 [21.36, 23.53] & \textbf{20.99} [20.12, 21.79] & yes \\
LibriSpeech-TTS & 33.54 [32.43, 34.62] & 33.74 [32.49, 34.88] & \cellcolor{red!10}no \\
\bottomrule
\end{tabular}
\end{table}

\newpage
\section*{Gain decomposition}
Splits the total EER change into two steps: $\Delta_1$ = effect of adding
self-training alone, $\Delta_2$ = effect of then adding consistency on top.
$\Delta_2$ helps on 3 targets but is $+0.35$ on LibriSpeech-TTS --- the one
target whose source AUC (.727) is far below the others ($\geq$ .84). Weak
source ranking gives unreliable pseudo-labels, and the consistency term then
amplifies that unreliability.
\begin{table}[h]
\centering
\begin{tabular}{@{}lcccc@{}}
\toprule
Target & Source EER & $\Delta_1$ (self-train) & $\Delta_2$ (consistency) & Total \\
\midrule
ASVspoof2019 & 5.03 & $-0.83$ & $-0.73$ & $-1.56$ \\
In-the-Wild & 12.78 & $-1.23$ & $-0.22$ & $-1.45$ \\
Arabic & 22.50 & $-0.30$ & $-1.20$ & $-1.50$ \\
LibriSpeech-TTS & 33.54 & $-0.15$ & \cellcolor{red!10}$+0.35$ & $+0.20$ \\
\bottomrule
\end{tabular}
\end{table}

\section*{Backbone comparison}
XLS-R = our SSL-pretrained encoder. RawNet2Lite = same epochs, same data, but
trained from random weights (no SSL pretraining). RawNet2Lite is at chance on
every target, so SSL pretraining --- not TTA --- is what makes a source model
cross-corpus viable in the first place. TTA operates on top of that.
\begin{table}[h]
\centering
\begin{tabular}{@{}lcc@{}}
\toprule
Target & XLS-R (ours) & RawNet2Lite (no SSL) \\
\midrule
ASVspoof2019 & 5.03 / .989 & 48.62 / .523 \\
LibriSpeech-TTS & 33.54 / .727 & 51.68 / .468 \\
In-the-Wild & 12.78 / .945 & 39.46 / .650 \\
Arabic & 22.50 / .843 & 56.52 / .416 \\
\bottomrule
\end{tabular}
\end{table}

\newpage
\section*{Comparison to train-time domain generalization (DANN, ASDG)}
Both baselines see the target's unlabeled audio during training --- more access
than our method gets --- and both are now at the same 10 seeds as ours, so the
head-to-head is no longer asymmetric. DANN is worse than Source-only on 3 of 4
targets and is the least seed-stable method tested. Ours beats both on every
target at roughly a third of DANN's compute (6.9 vs.\ 20.1 min per
target-seed).
\begin{table}[h]
\centering
\begin{tabular}{@{}lcccc@{}}
\toprule
Target & Source & DANN (20.1 min) & ASDG (8.2 min) & Ours (6.9 min) \\
\midrule
ASVspoof2019 & 5.03 / .989 & 7.35 / .966 & 4.98 / .987 & \textbf{3.47} / .994 \\
LibriSpeech-TTS & 33.54 / .727 & 34.43 / .717 & 32.64 / .749 & \textbf{33.74} / .727 \\
In-the-Wild & 12.78 / .945 & 13.35 / .924 & 17.33 / .891 & \textbf{11.33} / .958 \\
Arabic & 22.50 / .843 & 28.53 / .759 & 27.41 / .798 & \textbf{20.99} / .866 \\
\bottomrule
\end{tabular}
\end{table}

ASDG is the exception on LibriSpeech-TTS (32.64 vs.\ our 33.74) --- the same
weak-source-ranking target where our consistency term is a net negative.

\section*{Inductive check}
Rules out ``the model just memorised the scored clips'': adapt on one half of
the target, score on the disjoint other half. The main table's pattern
replicates (wins on 3 of 4), so the gain is real, not memorisation.
\begin{table}[h]
\centering
\begin{tabular}{@{}lccc@{}}
\toprule
Target & Source (half) & Ours (disjoint half) & Transductive gain, for reference \\
\midrule
ASVspoof2019 & 5.07 & \textbf{3.70} & $-1.56$ \\
In-the-Wild & 12.73 & \textbf{11.61} & $-1.45$ \\
Arabic & 22.31 & \textbf{21.47} & $-1.50$ \\
LibriSpeech-TTS & 33.45 & 33.58 & $+0.20$ \\
\bottomrule
\end{tabular}
\end{table}

\newpage
\section*{Our TTA on third-party public checkpoints (In-the-Wild)}
Every table above compares our source model against our baselines on our
protocol. This one runs the \emph{unmodified} published TTA config on weights
released by other groups, so the contribution stops depending on our own model
being the only comparator. Both checkpoint families are wav2vec2/XLS-R-300M +
AASIST; neither has seen In-the-Wild (31{,}779 clips). Only \textbf{16{,}706
parameters} are adapted (0.005\% of 315.9M): the top-4 LayerNorm blocks plus
the output layer.
\begin{table}[h]
\centering
\begin{tabular}{@{}llccc@{}}
\toprule
Checkpoint & Trained on & Source & Ours & $\Delta$ EER \\
\midrule
\texttt{ash56/ssl-aasist} & WaveFake / LJSpeech & 5.04 / .986 & \textbf{4.06} / .990 & $-0.98$ \\
DeepFense seed 42 & ASVspoof2019-LA & 16.41 / .918 & \textbf{10.10} / .961 & \cellcolor{green!12}$-6.31$ \\
DeepFense seed 240 & ASVspoof2019-LA & 8.41 / .970 & \textbf{7.88} / .973 & $-0.52$ \\
DeepFense seed 2 & ASVspoof2019-LA & 9.59 / .966 & \textbf{9.06} / .970 & $-0.53$ \\
\bottomrule
\end{tabular}
\end{table}

Two independent training corpora was deliberate: a gain that reproduces on both
is not a quirk of one model's score distribution. The best absolute detector is
\texttt{ash56}, and its gain is stable across seeds (4.06 / 4.09 / 4.09 on seeds
0/1/2). The largest gain goes to the \emph{weakest} source model, DeepFense
seed 42, which is what the ranking-vs-calibration account predicts.

\subsection*{Ablation and Tent on the same public weights}
\begin{table}[h]
\centering
\begin{tabular}{@{}lcccc@{}}
\toprule
Checkpoint & Source & Self-train only & Consistency only & Ours (both) \\
\midrule
\texttt{ash56/ssl-aasist} & 5.04 & 4.93 & 4.35 & \textbf{4.06} \\
DeepFense seed 42 & 16.41 & 11.99 & \cellcolor{red!10}28.65 & \textbf{10.10} \\
\bottomrule
\end{tabular}
\end{table}

On the strong checkpoint both components help and compose. On the weak one,
consistency \emph{alone} is actively harmful ($16.41 \to 28.65$) and
self-training is what rescues it --- the same pattern already documented on
LibriSpeech-TTS, now reproduced independently on third-party weights.

On DeepFense seed 42, \textbf{Tent collapses to 49.80\% EER / .502 AUC (exact
chance)} while ours takes the same checkpoint from 16.41 to 10.10. On a model
we did not train, with data we did not choose. This is the strongest evidence
that the \emph{structure} of the method, not our particular source model, is
doing the work.

\subsection*{Inductive check on public checkpoints}
\begin{table}[h]
\centering
\begin{tabular}{@{}lcc@{}}
\toprule
Checkpoint & Source (half) & Ours (disjoint half) \\
\midrule
\texttt{ash56/ssl-aasist} & 4.84 & \textbf{4.08} \\
DeepFense seed 42 & 15.07 & \textbf{10.21} \\
\bottomrule
\end{tabular}
\end{table}

\textbf{What is deliberately not claimed.} Neither checkpoint has a published
\emph{single-model} In-the-Wild EER (Garg et al.\ report 5.70\% for an
\emph{ensemble}; DeepFense ships no model card), so no claim of reproducing a
published figure is made. Correctness of the port rests instead on four gates,
all passed on CPU before any GPU time: strict weight load (421 tensors, 0
missing / 0 unexpected in the back-end); a check that the loaded weights differ
from a fresh XLS-R init; a score-polarity check (both public checkpoints use
$[\text{spoof}=0,\text{bonafide}=1]$, the opposite of our convention --- without
this check the whole experiment reports $100-x$); and an assert that the
trainable-parameter set is non-empty, since our selector hardcodes our own
module path and would otherwise silently adapt nothing.

\newpage
\section*{Domain divergence before/after adaptation}
$\hat d_\mathcal{A}$ and domain-probe accuracy both measure how easily a linear
classifier can tell source from target embeddings apart; both go \emph{up}
slightly after adaptation on every target. So adaptation works by re-shaping the
decision boundary, not by erasing domain identity --- which is exactly what DANN
optimizes for, and post-hoc explains why DANN underperforms here.
\begin{table}[h]
\centering
\begin{tabular}{@{}lcccc@{}}
\toprule
& \multicolumn{2}{c}{$\hat d_{\mathcal{A}}$} & \multicolumn{2}{c}{Domain probe acc.\ (\%)} \\
Target & source & adapted & source & adapted \\
\midrule
ASVspoof2019 & 1.572 & 1.609 & 89.3 & 90.2 \\
LibriSpeech-TTS & 1.460 & 1.504 & 86.5 & 87.6 \\
In-the-Wild & 1.335 & 1.357 & 83.4 & 83.9 \\
Arabic & 1.545 & 1.606 & 88.6 & 90.2 \\
\bottomrule
\end{tabular}
\end{table}

Single seed per target; this is a mechanism diagnostic, not a headline number.

\section*{Per-language MLAAD fake-recall}
Fake-clip recall averaged over 32 languages seen in training vs.\ 6 languages
held out of every source pool entirely. The 0.7-point gap is negligible ---
unseen languages are detected almost as well as seen ones.
\begin{table}[h]
\centering
\begin{tabular}{@{}lc@{}}
\toprule
Language group & Mean fake-clip recall \\
\midrule
Seen in source training (32 langs) & 93.9\% \\
Held out entirely (6 langs) & 93.3\% \\
\bottomrule
\end{tabular}
\end{table}

\section*{Known failure: official ASVspoof DF eval (Protocol A)}
The one protocol where the method fails outright, reported rather than omitted.
The target pool is heavily class-skewed, and our fixed symmetric $q=0.3$
pseudo-label split is the documented weak spot in that regime. Neither RawBoost
augmentation nor the adaptive arm rescues it; a prior-corrected variant recovers
part of the loss but does not reach source.
\begin{table}[h]
\centering
\begin{tabular}{@{}lcc@{}}
\toprule
Method & EER & AUC \\
\midrule
Source & \textbf{26.33} & .870 \\
Ours & \cellcolor{red!10}42.46 & .686 \\
Ours + prior correction & 32.58 & .820 \\
Ours + RawBoost & 42.49 & .686 \\
Ours + adaptive $q$/$E$ + RawBoost & 42.49 & .686 \\
\bottomrule
\end{tabular}
\end{table}

Single run on the official eval list (400{,}435 scored, 133{,}493 missing from
the local copy). This is the clearest target for the adaptive-$q$ work once GPU
frees up.

\end{document}
